\documentclass[11pt,a4paper]{article}
\usepackage[utf8]{inputenc}
\usepackage[T1]{fontenc}
\usepackage[french]{babel}
\usepackage{geometry}
\usepackage{xcolor}
\usepackage{hyperref}
\usepackage{titlesec}

\geometry{margin=1in}
\definecolor{titleblue}{RGB}{0,51,102}

\hypersetup{
    colorlinks=true,
    linkcolor=titleblue,
    urlcolor=titleblue,
    pdftitle={Lettre de Motivation - Ismail AISSAOUI},
}

\begin{document}

\noindent \textbf{Ismail AISSAOUI} \\
Ingénieur d'État en Intelligence Artificielle (ESI-SBA) \\
Email: ismail.aissaoui@example.com \\
Téléphone: [Votre Numéro] \\
LinkedIn: [Votre Lien]

\bigskip

\begin{flushright}
    À l'attention de M. Laurent AFONSO \\
    \textbf{AIXIAL Group (Groupe Alten)} \\
    Recherche \& Innovation \\
    Boulogne-Billancourt, France \\
    \medskip
    4 mai 2026
\end{flushright}

\bigskip

\noindent \textbf{Objet : Candidature au doctorat CIFRE : Modélisation dynamique multi-échelle des altérations biologiques dans les maladies neurodégénératives}

\bigskip

Monsieur Afonso,

Ingénieur d'État en Intelligence Artificielle diplômé de l'ESI-SBA, je vous adresse avec un vif intérêt ma candidature pour la thèse CIFRE intitulée « Modélisation dynamique multi-échelle des altérations biologiques dans les maladies neurodégénératives » au sein d'Aixial Group. Ma démarche est motivée par la volonté de mettre mon expertise en architectures d'IA avancées et en jumeaux numériques au service de la recherche clinique et des neurosciences computationnelles.

Actuellement en fin de cycle ingénieur chez Sonatrach, je réalise ma thèse sur le développement de **Jumeaux Numériques informés par la physique** pour la surveillance de systèmes complexes. Bien que mon application actuelle soit industrielle, les verrous scientifiques que je traite sont en parfaite adéquation avec votre projet :
\begin{itemize}
    \item \textbf{Modélisation Temporelle} : J'utilise les architectures \textbf{Mamba-SSM} et \textbf{xLSTM} pour modéliser des trajectoires dynamiques complexes. Ces modèles de « State Space » sont, selon moi, les outils les plus prometteurs pour capturer l'évolution temporelle des altérations biologiques et le déclin cognitif.
    \item \textbf{Intelligence Relationnelle} : Je maîtrise les \textbf{Réseaux de Neurones sur Graphes (GNN)}, essentiels pour articuler les données multi-échelles (moléculaires, cellulaires et cliniques) au sein d'un modèle intégratif cohérent.
    \item \textbf{Explicabilité et Rigueur} : Mon travail sur les modèles hybrides (mécanistes/IA) garantit non seulement la précision prédictive mais aussi la formulation d'hypothèses testables, répondant directement à votre exigence de modèles auto-explicatifs.
\end{itemize}

Le format CIFRE m'attire particulièrement. Mon expérience au sein d'une structure comme Sonatrach m'a doté d'une rigueur professionnelle et d'une capacité à communiquer mes résultats à des équipes pluridisciplinaires, des qualités indispensables pour évoluer au sein d'Aixial Group et collaborer avec vos partenaires du secteur pharmaceutique.

Je suis convaincu que mon profil d'« architecte système » appliqué à l'IA peut apporter un regard neuf et des outils mathématiques puissants à vos recherches sur les maladies neurodégénératives.

Je vous remercie de l'attention portée à mon dossier et reste à votre entière disposition pour un entretien afin de discuter de la manière dont mes compétences peuvent contribuer à l'innovation chez Aixial Group.

Je vous prie d'agréer, Monsieur, l'expression de mes salutations distinguées.

\bigskip

\noindent Ismail AISSAOUI

\end{document}
